%%%%%%%%%%%%%%%%%%%%%%%%%%%%%%%%%
%
%       EXERCÍCIO
%
%%%%%%%%%%%%%%%%%%%%%%%%%%%%%%%%%

\ifdefstring{\atividade}{prova}{%
    \renewcommand{\valorquestao}{\ValorQ\ pontos}
}%

\begin{exercicioBanco}[\valorquestao]
Um levantamento estatístico, feito em determinada região do país, mostrou que jovens com idades entre 4 e 17 anos assistem à televisão, em média, durante 6 horas por dia ($\mu = 6$). A tabela a seguir apresenta outras estatísticas produzidas por esse levantamento:

\begin{center}
\begin{tabular}{|c|c|}
\hline
\textbf{Estatística} & \textbf{Tempo ($T$, em horas)} \\ \hline
1.º quartil ($Q_1$) & 2 \\ \hline
2.º quartil ($Q_2$) & 4 \\ \hline
3.º quartil ($Q_3$) & 8 \\ \hline
1.º decil ($D_1$) & 1 \\ \hline
9.º decil ($D_9$) & 10 \\ \hline
\end{tabular}
\end{center}

Com base estritamente nas informações fornecidas na tabela, determine:

\begin{enumerate}[(a)]
    \item O coeficiente quartílico de assimetria de Bowley e classifique a distribuição quanto à assimetria.
    \item O coeficiente de curtose baseado em quartis e decis e classifique a curva de frequência quanto à curtose.
\end{enumerate}
\end{exercicioBanco}
